%- block legend
\vspace{-0.25cm}
\begin{center}
    \includegraphics[width=0.28\textwidth]{lobo.png}
\end{center}

Na floresta vive um lobo selvagem que adora melancias. Sempre que está com fome,
ele parte na floresta em busca da sua fruta favorita.\\
A floresta tem $N$ pontos e $M$ trilhas, que
ligam um par de pontos cada. As trilhas são as únicas formas de se mover pela
floresta e, para ir um ponto a outro, o lobo pode passar por várias trilhas.\\
São dadas $Q$ consultas. Em cada consulta, são dadas as posições atuais do
lobo e de uma melancia. Determine se o lobo consegue ir até essa melancia e saciar
sua fome.

%- endblock

%- block input

A primeira linha contém dois inteiros $N$ e $M$
($\VAR{N.min} \leq N \leq \VAR{N.max | sci}$,
$\VAR{M.min} \leq M \leq \VAR{M.max | sci}$), o número de pontos e de trilhas.
Considere que pontos são numerados de $1$ até $N$.\\
As próximas $M$ linhas descrevem uma trilha cada. Cada trilha é dada por dois
inteiros $A$ e $B$ ($1 \leq A, B \leq N$), indicando os pontos os quais a trilha
conecta.\\
A próxima linha contém um inteiro $Q$ ($\VAR{Q.min} \leq Q \leq \VAR{Q.max |
sci}$), o número de consultas.
As próximas $Q$ linhas descrevem uma consulta cada. Cada consulta é dada por
dois inteiros $l$ e $m$ ($1 \leq l, m \leq N$), indicando o ponto onde está o
lobo e a melancia, respectivamente.
%- endblock

%- block output
Para cada consulta, imprima \texttt{sim} se o lobo pode
ir até a melancia, ou \texttt{nao} caso contrário.
%- endblock

%- block explanation_0
Na primeira consulta, o lobo consegue chegar até a melancia percorrendo, por
exemplo, a trilha $1 \leftrightarrow 5$ e depois a trilha $5 \leftrightarrow 7$.
%- endblock

%- block editorial
Este problema envolve conceitos da Teoria dos Grafos.

Fazer uma busca (em profundidade ou em largura) para cada consulta não será rápido o bastante, dado que
cada busca tem custo $O(N+M)$ e o custo total seria, portanto,
$O(Q\times(N+M))$.

Entretanto, note que, em um grafo não direcionado, existe um caminho de $l$ para
$m$ se e somente se $l$ e $m$ estão na mesma componente conexa.

Assim, pré-processe o grafo com uma única busca (em profundidade ou em largura) para identificar
a componente conexa a qual cada vértice pertence. Então, para cada consulta,
verifique se $l$ e $m$ estão na mesma componente, ou em componentes distintas.

Referência a busca em profundidade: \url{https://cp-algorithms.com/graph/depth-first-search.html}  

Referência a busca em largura: \url{https://cp-algorithms.com/graph/breadth-first-search.html}

Como esta verificação pode ser feita em $O(1)$ após o pré-processamento, a complexidade total se torna
$O(N+M+Q)$.
%- endblock
